% ============================================================
% KDMILE 2026 — STAB Framework Paper
% Deadline: 08/jun/2026 | Limite: 8 páginas
% Submissão: https://jems3.sbc.org.br/kdmile2026
% ============================================================

\documentclass{kdmile}

\usepackage{graphicx,url}
\usepackage[T1]{fontenc}
\usepackage[utf8]{inputenc}
\usepackage{booktabs}
\usepackage{multirow}
\usepackage{listings}
\usepackage{float}

% kdmile.cls é baseado em ACM Trans — define resumo e abstract manualmente
\newenvironment{abstracten}{%
  \vspace{4pt}\begin{center}\normalsize\bfseries Abstract\end{center}%
  \noindent\small\ignorespaces
}{\normalsize\vspace{4pt}}

\newenvironment{resumo}{%
  \vspace{4pt}\begin{center}\normalsize\bfseries Resumo\end{center}%
  \noindent\small\ignorespaces
}{\normalsize\vspace{4pt}}

\providecommand{\address}[1]{}
\renewcommand{\institutename}{}

% ── Cabeçalho ───────────────────────────────────────────────
\markboth{Nogueira et al.}{STAB: A Complexity-Ladder Framework for Evaluating Local LLMs}

% ── Título ──────────────────────────────────────────────────
\title{STAB: A Complexity-Ladder Framework for Evaluating Local LLMs as Autonomous Agents}

% ── Autores ─────────────────────────────────────────────────
\author{
  João Conrado V. Nogueira\\
  {\small Pesquisador Independente, Macapá, AP, Brasil}
}

% ============================================================
\begin{document}
\maketitle
% Corrige rodapé: cls hardcoda KDMiLe 2025; sobrescreve após \maketitle
\gdef\@runningfoot{\@journalName, KDMiLe 2026.}
\gdef\@firstfoot{\@journalName, KDMiLe 2026.}
% ============================================================

% ── Abstract (EN) ───────────────────────────────────────────
\begin{abstracten}
Running LLMs as autonomous agents on local hardware is harder than it looks.
%
This paper reports STAB: a four-level complexity ladder (S1$\rightarrow$S4)
combined with a 12-hour temporal stability battery, built to expose
failure modes in on-premise tool-calling agents.
%
We tested eight models from four families (Qwen, Gemma, Llama, Mistral)
on a single NVIDIA RTX~3060 (12~GB VRAM).
%
The central finding is that most failures trace back to the instruction and
contextual decision layer—not to raw model capacity.
%
Two phenomena stood out: \emph{tool-skipping} triggered by descriptive
filenames (H12), and the orthogonality between format compliance and semantic
accuracy (H10).
%
\texttt{gemma4:e4b} was the only model to pass all four levels without a
single failure.

\textbf{Keywords:} Local LLMs, Autonomous Agents, Tool Calling, Instruction
Following, Small Language Models, On-Premise Benchmark.
\end{abstracten}

% ── Resumo (PT) ─────────────────────────────────────────────
\begin{resumo}
O uso de Large Language Models (LLMs) em sistemas agênticos locais enfrenta
desafios críticos de consistência e confiabilidade, especialmente em hardware
de consumo.
%
Este trabalho apresenta o protocolo STAB, uma metodologia de avaliação em escada
de complexidade (S1$\rightarrow$S4) combinada a um protocolo de estabilidade
temporal para diagnosticar falhas em agentes on-premise.
%
Avaliamos oito modelos de quatro famílias (Qwen, Gemma, Llama e Mistral) em
servidor com GPU NVIDIA RTX~3060 (12~GB VRAM).
%
Os resultados revelam que as falhas agênticas originam-se predominantemente na
camada de instrução e decisão contextual, e não na capacidade bruta dos modelos.
%
Identificamos o fenômeno de \emph{tool-skipping} por nomeação semântica (H12) e
demonstramos que conformidade de formato e acurácia semântica são dimensões
ortogonais (H10).
%
O modelo \texttt{gemma4:e4b} destacou-se como referência de maior
confiabilidade, mantendo desempenho perfeito em todos os níveis de complexidade.

\textbf{Palavras-chave:} LLMs Locais, Agentes Autônomos, Tool Calling,
Instruction Following, Small Language Models, Benchmark On-Premise.
\end{resumo}

% ============================================================
\section{Introduction}
\label{sec:intro}

Getting a local LLM to reliably call tools—not just chat—is harder than it looks.
Models that score well on standard benchmarks often break in agentic settings:
they skip tool invocations, hallucinate results, or return correct JSON
wrapped in markdown that the calling script cannot parse.
Cloud-based evaluations such as BFCL and AgentBench do not surface these issues
because they assume virtually unlimited compute and allow multi-pass retries.
On a single 12\,GB GPU, $T=0$, and one-shot execution, the failure modes look
quite different—and no existing benchmark was designed specifically to expose them.

% ── Parágrafo 2: contribuições ──────────────────────────────
This paper introduces STAB to fill that gap. The main contributions are:

\begin{itemize}
  \item \textbf{STAB Framework:} A four-level methodology (S1$\rightarrow$S4)
        evaluating from basic tool use to autonomous multi-step planning,
        validated by a temporal stability protocol.

  \item \textbf{Failure Taxonomy:} Classification of agentic errors across four
        layers: protocol, decision, format, and grounding.

  \item \textbf{H12 — Semantic Tool-Skipping:} Empirical identification of
        tool invocation omission induced by descriptive file names.

  \item \textbf{H10 — Orthogonality:} Evidence that format compliance and
        semantic accuracy are independent evaluation dimensions.
\end{itemize}

% ── Parágrafo 3: organização ────────────────────────────────
The paper is organized as follows: Section~\ref{sec:related} reviews related
work; Section~\ref{sec:method} describes the STAB methodology;
Section~\ref{sec:results} presents results and analysis;
Section~\ref{sec:discussion} discusses findings; and
Section~\ref{sec:conclusion} concludes the work.

% ============================================================
\section{Related Work}
\label{sec:related}

\textbf{Agent benchmarks and tool use.}
The Berkeley Function Calling Leaderboard (BFCL)~\cite{patil2025} is the
de-facto benchmark for tool evaluation, focusing predominantly on cloud models
and large LLMs.
AgentBench~\cite{liu2024} proposes multi-environment evaluation across eight
domains, identifying instruction following as the primary obstacle to autonomy.
SLM-Bench~\cite{pham2025} introduces environmental impact metrics and
controlled hardware setups, approaching our local evaluation methodology.

\textbf{Failure modes and instruction following.}
A taxonomy of 14 failure modes in multi-agent systems is proposed
in~\cite{cemri2025}, highlighting specification disobedience.
SLM performance for function calling on edge devices is mapped
in~\cite{kavathekar2025}, while~\cite{jhandi2026} demonstrates that SLMs with
up to 350M parameters can outperform larger models in tool use with targeted
fine-tuning. A decline in instruction compliance in smaller models is documented
through the KCIF dataset in~\cite{murthy2024}, and training pipelines for
enterprise scenarios requiring deterministic tool use are explored
in~\cite{zeng2024}.

\textbf{Differentiation.}
STAB is positioned differently: it targets on-premise hardware, introduces
a temporal stability battery, and explicitly probes contextual biases such as
semantic naming—none of which appear in the benchmarks reviewed above.

% ============================================================
\section{Methodology}
\label{sec:method}

\subsection{Experimental Setup}

\begin{table}[H]
\centering
\caption{Experimental environment.}
\label{tab:setup}
\begin{tabular}{ll}
\toprule
\textbf{Component} & \textbf{Specification} \\
\midrule
CPU     & Intel Core i5-10400F (6c/12t) \\
RAM     & 32\,GB DDR4 \\
GPU     & NVIDIA RTX 3060 (12\,GB VRAM) \\
OS      & Ubuntu 24.04 \\
Runtime & Ollama (REST \texttt{/api/chat}), $T=0$, \texttt{stream=false} \\
\bottomrule
\end{tabular}
\end{table}

\subsection{STAB Framework — Complexity Ladder}

The STAB framework organizes evaluation into four complexity levels, each
requiring greater cognitive and operational capability from the agent:

\begin{table}[H]
\centering
\caption{STAB complexity ladder.}
\label{tab:stab}
\begin{tabular}{clllcc}
\toprule
\textbf{Level} & \textbf{Name} & \textbf{Tool} & \textbf{Objective} & \textbf{K} & \textbf{Key Metric} \\
\midrule
S1 & Protocol      & \texttt{run\_bash}  & Tool invocation + return     & 5 & TOOL-ACC \\
S2 & Extraction    & \texttt{run\_bash}  & Structured JSON output       & 5 & JSON-ACC \\
S3 & Log Analysis  & \texttt{read\_file} & Semantic classification      & 4 & CLASS-ACC \\
S4 & Multi-log     & \texttt{read\_file} & Multi-step triage \& ranking & 3 & TRIAGE-ACC \\
\bottomrule
\end{tabular}
\end{table}

The framework utilizes synthetic fixtures rather than production logs to ensure
reproducibility and variable control. Each fixture consists of an input file
(.log) and a ground truth JSON file (.expected.json). The ground truth employs a
controlled vocabulary for issue identification (e.g., \texttt{high\_cpu\_usage},
\texttt{database\_unavailable}), transforming open-ended evaluation into
rigorous classification and extraction.

Each test follows a two-turn execution protocol via the Ollama REST API
(\texttt{/api/chat}). In Turn 1, the model receives the prompt and the JSON
Schema for the tools, producing a structured \texttt{tool\_call}. The script
executes the operation on the host and captures \texttt{stdout}. In Turn 2, this
result is returned to the model, which generates the final JSON response.
Determinism is ensured by setting temperature to zero and \texttt{stream=false}.

\subsection{Models Evaluated}

\begin{table}[H]
\centering
\caption{Models evaluated.}
\label{tab:models}
\begin{tabular}{llrll}
\toprule
\textbf{Model} & \textbf{Family} & \textbf{VRAM (MiB)} & \textbf{Type} & \textbf{Final Status} \\
\midrule
\texttt{gemma4:e4b}             & Gemma4    & 9{,}821 & Zero-shot  & \textbf{Reference} \\
\texttt{qwen3.5:9b}             & Qwen3.5   & 7{,}953 & Zero-shot  & Approved (S2B) \\
\texttt{qwen3:14b}              & Qwen3     & 9{,}437 & Zero-shot  & Approved (S1--S3) \\
\texttt{llama3.1:8b}            & Llama3.1  & 4{,}868 & Zero-shot  & Approved (Partial) \\
\texttt{llama3-groq-tool-use:8b}& Llama3    & 4{,}754 & Fine-tuned & Approved (Partial) \\
\texttt{mistral:7b}             & Mistral   & 4{,}458 & Zero-shot  & Discarded (H14) \\
\texttt{phi4-mini}              & Phi4      & 2{,}500 & Zero-shot  & Discarded (H13) \\
\texttt{qwen3.5:2b}             & Qwen3.5   & 3{,}829 & Zero-shot  & Discarded (H16) \\
\bottomrule
\end{tabular}
\end{table}

\subsection{Metrics}

Evaluation employs an ordinal scale from 0 to 4, where 4 represents a perfect
run (correct tool invocation, valid JSON output, and accurate semantic content)
and 0 represents complete failure. Metrics are grouped into three categories:
consistency metrics (CON, TOOL-CON) assess run-to-run stability; instruction and
accuracy metrics (INST, TOOL-ACC, JSON-ACC, CLASS-ACC, PARSE-ACC, CHAIN-ACC,
TRIAGE-ACC) measure correctness at each STAB level; and the composite STAB score
reflects temporal stability across the full 12-hour battery protocol.

\begin{table}[H]
\centering
\caption{Evaluation metrics (ordinal scale 0--4).}
\label{tab:metrics}
\begin{tabular}{ll}
\toprule
\textbf{Metric} & \textbf{Description} \\
\midrule
CON       & Consistency: structural and decision stability across runs \\
TOOL-CON  & Tool consistency: repeated valid tool calls \\
INST      & Instruction adherence: format and constraint compliance \\
TOOL-ACC  & Functional correctness: right tool, successful execution \\
JSON-ACC  & JSON accuracy: schema validity and value precision \\
CLASS-ACC & Classification accuracy: correct global severity \\
PARSE-ACC & Extraction accuracy: host, issue, timestamp precision \\
CHAIN-ACC & Chaining accuracy: correct number of tool calls (e.g. 3/3) \\
TRIAGE-ACC& Triage accuracy: correct priority rank and escalation rule \\
STAB      & Temporal stability: score maintenance across 12\,h batteries \\
\bottomrule
\end{tabular}
\end{table}

\subsection{Failure Taxonomy}

Failures were classified across five layers based on their root cause.
Protocol-layer failures involve structural violations of the tool-calling format.
Decision-layer failures indicate contextual reasoning errors, such as inferring
results without executing the required tool. Format-layer failures produce
syntactically correct but non-compliant output. Grounding failures generate
structurally valid but factually fabricated responses. Instruction-layer failures
stem from ambiguous action interpretation. Each failure is assigned a unique
code: H-prefixed codes indicate hypothesis-driven findings confirmed through
controlled experiments, while AP codes mark anomalous patterns identified during
standard execution runs.

\begin{table}[H]
\centering
\caption{Identified failure taxonomy.}
\label{tab:failures}
\begin{tabular}{lllll}
\toprule
\textbf{Layer} & \textbf{Type} & \textbf{Code} & \textbf{Concrete Example} & \textbf{Model} \\
\midrule
Protocol  & Protocol Hallucination  & H13 & JSON in \texttt{content} instead of \texttt{tool\_calls} & phi4-mini \\
Protocol  & Echo-tool               & H16 & Delivers JSON via \texttt{run\_bash echo}                & qwen3.5:2b \\
Decision  & Semantic Tool-Skipping  & H12 & Skipped read; inferred severity from filename            & qwen3:14b \\
Format    & Markdown Wrapper        & AP2 & Correct JSON wrapped in \texttt{```json}                 & qwen3.5:9b \\
Grounding & Silent Grounding        & H14 & Perfect JSON with fabricated hardware data               & mistral:7b \\
Instruction& Action Ambiguity       & AP1 & Used \texttt{cat} instead of executing script            & qwen3.5:9b \\
\bottomrule
\end{tabular}
\end{table}

% ============================================================
\section{Results}
\label{sec:results}

\subsection{Overall Comparison}

Three models failed to clear the early levels: \texttt{phi4-mini} (H13, S1),
\texttt{qwen3.5:2b} (H16, S2), and \texttt{mistral:7b} (H14, S2)—each for a
different structural reason.
Of the five models that continued, \texttt{gemma4:e4b} was the only one to pass
every level cleanly (S1$\rightarrow$S4; see Table~\ref{tab:results}).
The sharpest drop appeared at S4: models that had handled S1--S3 without issues
collapsed on multi-step triage, pointing to a decision-layer ceiling that
single-step tasks simply do not reveal.

\begin{table}[H]
\centering
\caption{STAB scores per model and level (0--4; \textbf{bold} = perfect run).}
\label{tab:results}
\begin{tabular}{lccccccl}
\toprule
\textbf{Model} & \textbf{S1} & \textbf{S2} & \textbf{S3} & \textbf{S4-A} & \textbf{S4-B} & \textbf{S4-C} & \textbf{S1→S4} \\
\midrule
\texttt{gemma4:e4b}              & \textbf{4} & \textbf{4} & \textbf{4} & \textbf{4} & \textbf{4} & \textbf{4} & \textbf{Yes} \\
\texttt{qwen3.5:9b}              & 4 & 4 (S2B) & 2$^a$ & 4 & 2 & 2 & No \\
\texttt{qwen3:14b}               & 4 & 4 & 4 & 0 & 0 & 0 & No \\
\texttt{llama3.1:8b}             & 4 & 4 & 2 & 4 & 2 & 3 & No \\
\texttt{llama3-groq-tool-use:8b} & 4 & 4 & 4 & 4 & 2 & 2 & No \\
\texttt{qwen3.5:2b}              & 4 & disc. H16 & --- & --- & --- & --- & No \\
\texttt{mistral:7b}              & 4 & disc. H14 & --- & --- & --- & --- & No \\
\texttt{phi4-mini}               & disc. H13 & --- & --- & --- & --- & --- & No \\
\bottomrule
\end{tabular}
\begin{minipage}{\linewidth}
\vspace{2pt}
\footnotesize
$^a$\,\texttt{qwen3.5:9b} achieved CLASS-ACC=4 in S3-mixed but CON=INST=2 due to
persistent markdown wrapper, illustrating the orthogonality of dimensions (H10).
\end{minipage}
\end{table}

\subsection{Key Findings}

\textbf{H12 — Semantic Tool-Skipping.}
The \texttt{qwen3:14b} model skipped file reading in S4 scenarios when filenames
were semantically descriptive (e.g., \texttt{critical-001.log}), inferring
severity without executing the tool. In the controlled H12 experiment with
opaque filenames, the model correctly invoked the tool, confirming the
contextual bias hypothesis. \textbf{Mitigation:} use opaque filenames or
explicit mandatory-read instructions.

\textbf{H10 — Orthogonality of Format and Semantic Accuracy.}
Format compliance and semantic accuracy proved to be independent dimensions.
A model can produce semantically perfect output wrapped in an invalid format
(markdown), or a structurally valid JSON containing hallucinated data.
This was exemplified by \texttt{qwen3.5:9b}, which maintained high semantic
accuracy while failing formatting constraints.

\textbf{gemma4:e4b as Reference Model.}
The \texttt{gemma4:e4b} model was the only one to achieve perfect scores across
all levels (S1$\rightarrow$S4) with STAB=4 and a throughput of 73.0\,tok/s.
Its internal reasoning (\emph{thinking}) activated selectively for high-complexity
disambiguation scenarios (S4-C, such as timestamp-based tie-breaking),
indicating adaptive cognitive scaling.

\subsection{Throughput and Latency}

\begin{figure}[H]
  \centering
  \includegraphics[width=\linewidth]{plot-throughput-kdmile.png}
  \caption{Average throughput per model and scenario (tok/s), Battery 1.
           Run~1 excluded (cold start). Variation $<0.2\%$ across batteries
           confirms hardware-bound throughput.}
  \label{fig:throughput}
\end{figure}

Variation between batteries stayed below 0.2\% for all models, confirming
throughput is hardware-bound rather than prompt-dependent.
Run~1 of each battery was excluded from the averages (cold start penalty).
Three performance clusters emerged: \texttt{gemma4:e4b} and the Llama family
at the top, \texttt{qwen3.5:9b} in the middle, and \texttt{qwen3:14b} at the
bottom—its S2 latency was inflated by a large internal reasoning trace.

\subsection{Temporal Stability}

All five approved models completed Battery~2 without degradation: agentic
decisions and semantic accuracy held at STAB=4 throughout.
The one anomaly was \texttt{qwen3:14b}, whose internal reasoning volume varied
by 7\% between batteries—but the final output was identical each time, so this
is a structural quirk of the model's chain-of-thought, not actual instability.

% ============================================================
\section{Discussion}
\label{sec:discussion}

\textbf{Failures are instructional, not capacity-based.}
Two corrective experiments in this study tell the same story.
Adding an explicit format instruction recovered \texttt{qwen3.5:9b}'s S2B score
from 2 to 4. Adding a mandatory-read directive eliminated \texttt{qwen3:14b}'s
tool-skipping entirely.
Neither fix required changing the model or retraining—only the prompt.
Model size is therefore a poor proxy for agentic reliability: models above 8B
clearly have the underlying logic; what breaks them is context misalignment,
and that is fixable at the instruction layer.

\textbf{Fine-tuned vs.\ zero-shot in on-premise environments.}
The head-to-head between \path{llama3-groq-tool-use:8b} (fine-tuned) and
\texttt{llama3.1:8b} (zero-shot) produced a split result.
On S3-mixed, which requires applying binary precedence rules, the fine-tuned
model won clearly (Groq=4 vs.\ Llama=2)—specialization works for single-step
structural compliance.
But on S4-B and S4-C, both failed.
This raises a concrete question for the field: does tool-use fine-tuning improve
structural adherence without touching the multi-step reasoning needed for chained
decisions? The data here says yes—and that is a meaningful ceiling in current
tuning pipelines.

\textbf{Practical recommendations.}
On hardware with 12\,GB VRAM, \texttt{gemma4:e4b} is the clear choice: perfect
consistency, 73\,tok/s, no prompt engineering required.
Below 6\,GB, Llama3.1 works—but needs a thin post-processing layer to handle
format failures.
The H12 finding also has a security angle: descriptive filenames allow models
to infer answers and skip the tool call entirely, which means the tool is never
actually executed.
Using opaque identifiers, or appending \texttt{[READ REQUIRED]} to the instruction,
closes that gap.

\textbf{Limitations.}
The study ran on a single hardware configuration with synthetic fixtures only.
Those are real constraints, but they also made full reproducibility possible
and kept variables isolated—a deliberate trade-off.
The most operationally significant result to carry forward is H14 (silent
grounding in Mistral): structurally perfect output that is factually wrong
will not be caught by a JSON parser.
External semantic validators are not optional in production—they are the only
safeguard against that specific failure mode.

% ============================================================
\section{Conclusion}
\label{sec:conclusion}

Local agent reliability is a three-variable problem: model robustness,
instruction clarity, and context neutrality.
All three must hold simultaneously—failure in any one of them breaks the chain.
\texttt{gemma4:e4b} was the only model in this study to keep all three in check
across every level, which shows that quantized MoE architectures have crossed a
practical threshold for structured agentic work on consumer hardware.

The failures in 8B--14B models were, almost without exception, fixable through
prompt engineering alone—no retraining needed.
That is an encouraging result for anyone building agents under hardware constraints.

What this study did not cover: real production logs, diverse hardware profiles,
and grammar-constrained decoding as an alternative to prompt-based format
enforcement.
The 12-hour temporal battery also leaves open whether longer deployments surface
degradation patterns that were not visible here.
Those are the natural extensions—not by convention, but because the data from
this study raises those specific questions.

% ============================================================
\bibliographystyle{kdmile}
\bibliography{references}

\end{document}
